\documentclass[12pt]{article} 
\usepackage[margin=0.5in]{geometry} 
\usepackage{amsmath,amssymb} 
\usepackage{graphicx} 
\usepackage{tikz} 
\usepackage{natbib}
\usepackage{hyperref}

\bibliographystyle{apalike}
\title{Echo Chambers and Diagnostic Expectations: A Behavioral Macroeconomics Research Proposal} \author{Second-Year Ph.D. Proposal} \date{\today} 

\begin{document}
\maketitle

\section{Introduction}

\paragraph{Project Summary.} \emph{This project investigates how exposure to polarized narrative ``echo chambers''---measured via podcast media consumption---affects diagnostic overreaction in macroeconomic expectation updating, with heterogeneity by socio-economic status (SES). We propose a shift--share (Bartik) identification strategy combined with microdata on expectations to estimate how narrative-driven signal distortion (captured by a diagnostic expectation parameter $\theta$) varies with media exposure and SES, and how this distortion influences household decisions relevant to inequality.}

A central simplification in macro is that households form beliefs as if they observed (or could costlessly infer) the relevant aggregate state: the full-information rational expectations (FIRE) benchmark. In practice, households learn about inflation, labor-market conditions, and policy largely through \emph{information intermediaries}---news, social media, and increasingly long-form audio---and these intermediaries are not neutral aggregators. Modern media markets reward attention and engagement, and audience segmentation creates strong incentives to supply one-sided frames that confirm priors and sustain loyal followings.

A key open question---and the focus of this proposal---is whether the content and structure of information sources can causally drive these expectation biases. In particular, the rise of \textbf{polarized narratives in media} (e.g. partisan podcasts and online news) may inject \emph{correlated noise} into agents' information sets, leading to shared belief distortions or ``echo chambers.'' If a large subset of households receives information from the same biased narrative source, they may update their beliefs in a similarly distorted way, causing systematic forecast errors (e.g. over-optimism or over-pessimism) that violate rational expectations. Moreover, different socio-economic groups may have differential exposure to such narratives (due to differences in media consumption habits or trust) and potentially different susceptibility to bias, leading to \emph{heterogeneity} in expectation errors across SES groups.

At the individual level, there is evidence of \emph{over}-reaction: individual forecasts overweight new signals, leading to a negative correlation between forecast errors and prior revisions and this is consistent with the \textit{diagnostic expectations} model of \citet{BordaloGennaioliMaShleifer2020} (henceforth BGMS), in which agents exaggerate the importance of recent signals. These biases in expectation formation are important because household expectations about inflation, income, house prices, etc., influence their consumption, saving, and investment decisions, thereby impacting macroeconomic dynamics and wealth inequality.

Models of selective memory and cue-triggered retrieval further imply that news and contextual cues can shift beliefs even when they carry little new statistical information, by bringing particular experiences or associations to mind and crowding out others \citep{BordaloGennaioliShleifer2020MemoryAttentionChoice}. Recent evidence links these mechanisms directly to persuasion and narratives: narratives or political advertising can change beliefs by activating otherwise neglected experiences, generating bias and heterogeneity even under common public information \citep{BordaloConlonGennaioliKwonShleifer2026HowPeopleUseStatistics,BordaloBurroCoffmanGennaioliShleifer2025ImaginingFuture}.

These mechanisms naturally map into modern media environments. Polarized podcasts and narrative-driven outlets may act as \emph{attention and memory cues} that systematically emphasize particular framings (e.g., ``inflation is out of control'', ``the economy is rigged'', ``soft landing is assured''), thereby distorting the effective signal that households extract from macro news. Because such narratives are \emph{shared} within media communities, they can generate \emph{correlated} belief distortions across exposed households (an ``echo chamber'' channel), rather than purely idiosyncratic noise. Moreover, exposure and susceptibility may vary with SES, creating systematic heterogeneity in forecast errors and downstream choices. Since belief distortions can translate into real decisions with distributional consequences (e.g., reluctance to realize nominal losses in household balance-sheet decisions), understanding narrative-driven expectation biases is macroeconomically and distributionally relevant \citep{GarnacheSchmalzSommervollStickyAssetPrices}.

Against this backdrop, survey evidence rejects FIRE in a disciplined way: household expectations exhibit systematic \emph{information stickiness} and heterogeneity in updating frequency. In particular, \citet{difrancesco2024_sticky_information_wealth} documents (using survey data) that households update aggregate expectations infrequently and that \emph{wealthier households update more frequently}, and shows that this heterogeneity matters quantitatively for monetary transmission in a HANK environment.

This provides a natural benchmark for our heterogeneity focus: SES (income/education/wealth) shapes either exposure to information or the ability/willingness to process it, so expectation dynamics and policy transmission may be inherently distributional even before we add narrative polarization.

Crucially, infact, deviations from FIRE need not be only about \emph{infrequent} updating; they can also be about \emph{distorted} updating when the information environment is systematically slanted. Behavioral models such as diagnostic expectations formalize over-weighting of salient or one-sided signals, generating predictable forecast errors from forecast revisions through a diagnosticity wedge $\theta$.


The motivation for focusing on podcasts as a proxy for narrative exposure is twofold. First, podcasts (and similar digital media) have become an increasingly important channel by which Americans consume news and economic commentary. Over half of U.S. adults (54\%) listened to a podcast in the past year \citep{pew_podcasts_news_fact_sheet_2025}, and about one-third (32\%) report that they sometimes get news from podcasts \citep{pew_podcasts_news_fact_sheet_2025}. Political and economic commentary podcasts often present highly curated, opinionated interpretations of economic events (for example, portraying jobs reports or inflation data through an ideologically charged lens). Such content can create polarized “storylines” about the economy. Second, podcasts allow measurement: we can scrape metadata and transcripts from sources like \textit{Listen Notes} and identify narrative sentiment or slant at the show or episode level. This makes it feasible to construct quantitative measures of \textbf{echo chamber exposure} for different regions or groups.


\paragraph{Behavioral Macro Framing and Contribution.} Our study contributes to behavioral macroeconomics by examining expectations through the joint lens of \emph{diagnostic expectations} theory \citep{BordaloGennaioliMaShleifer2020} and \emph{narrative economics} \citep{Shiller2017}. Diagnostic expectations posit that agents "overreact" to new information, placing excessive weight on signals that are representative of recent developments, which can generate momentum-like forecast errors \citep{GEMMI2026103868}. Narrative economics emphasizes that popular stories or narratives can drive economic behavior. We integrate these ideas: biased narratives in media serve as correlated signals that may trigger diagnostic overreaction. Unlike idiosyncratic errors, narrative-driven expectation errors could be correlated across agents, potentially leading to aggregate effects (e.g. self-fulfilling pessimism booms or busts) and distributional consequences (if certain groups are more influenced by such narratives).
Our project also speaks to the literature on information frictions \citep[e.g.][]{CoibionGorodnichenko2015} and rational inattention \citep{sims2003_implications_rational_inattention}. Whereas rational inattention models explain some forecast sluggishness by cognitive processing limits \citep{sims2003_implications_rational_inattention}, they typically do not generate systematic overreaction of the kind BGMS document. By focusing on \textbf{echo chambers}, we explore an alternative mechanism: agents optimally but naively trust their information sources, not realizing that these sources may be biased or incomplete (i.e., they treat narrative signals as truth, triggering diagnostic over-updating). This can be viewed as a form of \textit{endogenous noise} in information processing, as opposed to the exogenous noise in standard noisy-signal models \citep{sims2003_implications_rational_inattention}. In essence, exposure to an echo chamber could induce what we term a \textbf{narrative-driven diagnostic distortion} in expectations.


\paragraph{Research Questions.} The main research question driving this proposal is:
\begin{quote}
    \textbf{Does exposure to polarized narrative “echo chambers” causally induce diagnostic overreaction in household macroeconomic expectations, and if so, how does this effect differ across socio-economic groups?}
\end{quote}
From this, we derive several secondary questions:
\begin{itemize}
    \item \textbf{Q1a:} What are the measurable dimensions of narrative echo chamber exposure (e.g. intensity of partisan content, direction of slant, fragmentation of sources), and which of these most strongly influence expectation formation?
    \item \textbf{Q1b:} To what extent do households with different SES (e.g. education or income levels) differ in their baseline expectations errors and in their responsiveness to narrative exposure? In particular, are less-educated or lower-income households more prone to narrative-induced overreaction (or underreaction) than higher-SES households?
    \item \textbf{Q1c:} Do distortions in expectations linked to narrative exposure translate into differences in household \emph{decisions} on inequality-relevant margins (such as consumption spending, credit uptake, investment, or savings behavior)? In other words, are there tangible economic consequences of echo chamber-induced belief distortions that could affect the distribution of outcomes (wealth, financial stability, etc.) across households?
\end{itemize}


\paragraph{Hypotheses.} We will empirically test the following key hypotheses:
\begin{enumerate}
    \item \textbf{H1: Echo Chamber Overreaction.} Households in regions with greater exposure to polarized economic narratives (high ``echo chamber intensity'') exhibit \emph{diagnostic overreaction} in their expectation updating. In practice, this means their forecast revisions overshoot what a fully rational agent would do, leading to forecast errors that are \emph{predictably} in the opposite direction of their revisions \citep{GEMMI2026103868}. For example, after a positive economic news shock, high-exposure households revise expectations upward too strongly and subsequently overestimate actual outcomes (yielding a negative error).
    \item \textbf{H2: SES Heterogeneity in $\theta$.} Socio-economic status moderates the impact of narrative exposure on expectations. We hypothesize that lower-SES households (e.g. those with lower education or income) will have a larger diagnosed distortion parameter $\theta$ (stronger overweighting of narrative signals) compared to higher-SES households, given the same exposure. This could be due to differences in financial literacy or alternative information sources. In contrast, higher-SES individuals might filter or counteract biased narratives more effectively, exhibiting smaller $\theta$ (closer to rational updating).
    \item \textbf{H3: Directional Asymmetry.} The \emph{direction} of narrative slant matters for belief distortions. Households primarily exposed to pessimistic/negative economic narratives will bias their expectations downward (overly pessimistic forecasts), whereas those in the opposite echo chamber (optimistic or bullish narratives) will bias upward. This yields an \emph{expectation polarization}: for instance, right-leaning economic narratives might consistently predict higher inflation or stronger growth than left-leaning narratives, leading their respective audiences to hold divergent expectations.
    \item \textbf{H4: Fragmentation and Disagreement.} More \textbf{fragmented} media environments (where the audience is split among many echo chambers rather than one dominant information source) lead to greater \emph{disagreement} in expectations across individuals. In regions with high narrative fragmentation, we expect to observe higher cross-sectional dispersion in forecasted inflation or growth. By contrast, in regions where one narrative dominates, households' expectations may be more biased on average but also more homogeneous within that area.
    \item \textbf{H5: Real Effects on Behavior and Inequality.} If echo chamber exposure indeed distorts expectations, these distorted beliefs will influence behavior in ways that could affect inequality. For example, households with over-optimistic income expectations might over-borrow or under-save, risking future financial stress, whereas overly pessimistic households might under-invest in assets (like not buying stocks or a home due to dire economic outlook). We hypothesize that areas/groups with higher narrative-induced $\theta$ will show measurable differences in outcomes such as consumption growth, credit default rates, or investment portfolio allocations. Over time, these behavioral responses could contribute to widening gaps (e.g. if one group consistently misinterprets the economy and makes suboptimal financial decisions, their wealth may grow more slowly relative to a less-biased group).
\end{enumerate}
By testing these hypotheses, the project aims to shed light on the causal chain from media narratives to belief formation to economic actions. Next, we discuss some motivating evidence and facts that inform these hypotheses and motivate the empirical approach.


\section{Motivating Evidence and Stylized Facts}
A number of empirical findings motivate our study, illustrating both the existence of biases in expectations and the potential role of narratives and heterogeneity:

\subsection*{Biases and Predictable Errors in Expectations}
Surveys of expectations have revealed patterns inconsistent with the textbook rational expectation hypothesis. One widely documented fact is that forecast errors are often \emph{predictable by past forecast revisions}, violating the no-predictability condition of rational expectations \citep{GEMMI2026103868}. In particular, \citet{CoibionGorodnichenko2015} show that when forecasters update their predictions, they tend to do so insufficiently on average: if the consensus forecast is revised upward, subsequent actual outcomes are typically under-predicted (positive forecast errors) — evidence of \emph{under-reaction} or information rigidity (sometimes termed ``sticky expectations'') \citep{GEMMI2026103868}. This implies agents incorporate news only partially, consistent with models of imperfect information or rational inattention.
On the other hand, recent work has identified \emph{over}-reaction at the individual forecaster level. \citet{BordaloGennaioliMaShleifer2020} (BGMS) document that individual forecasts in the Survey of Professional Forecasters exhibit \textit{excess sensitivity} to new information: an individual who sharply revises their forecast up tends to over-shoot, making an error in the opposite direction (actual outcome falls short of their forecast). In regression terms, there is a significantly \textbf{negative correlation between individual forecast errors and forecast revisions} \citep{GEMMI2026103868}, whereas rational expectations predict zero correlation. This is exactly the pattern diagnostic expectations would generate \citep{BordaloGennaioliMaShleifer2020}. The intuition is that an agent overweighs an incoming signal (e.g. an unexpectedly good GDP growth indicator) relative to prior baseline, leading them to forecast too high. Subsequent reality, influenced by many factors, is on average more moderate, so a negative error results. This phenomenon has been observed not just for professional forecasters but also in households and firms to some extent, especially when news shocks are salient.
These two facts — aggregate underreaction and individual overreaction — can coexist and even stem from a common behavioral model once we consider heterogeneous agents and strategic or noisy reporting \citep{BordaloGennaioliMaShleifer2020, AngeletosLian2018}. For our purposes, they establish that expectation formation is \emph{behavioral} (deviating from FIRE) in ways that need explanation. One plausible explanation for overreaction is precisely the presence of \textbf{diagnostic narratives}: agents might latch onto vivid stories or signals and overweight them in their belief updating.

\subsection*{SES Differences in Expectations and Biases}
Socio-economic differences in expectations have been widely noted in survey data. Individuals with lower income or education often have less accurate and less ``anchored'' expectations, especially about macroeconomic variables like inflation. For instance, lower-income households in the NY Fed Survey of Consumer Expectations (SCE) report higher inflation expectations and are more sensitive to short-term price shocks than higher-income households. Recent evidence shows that the upper quantile of low-income respondents has inflation expectations that are not well anchored to the central bank's target, and they react more strongly to monetary policy shocks \citep{doh_lee_park_2024_rwp2406}. In other words, when a surprise rate hike occurs, less wealthy households adjust their expectations more dramatically (perhaps indicating a lower signal-to-noise ratio in their information or a tendency to overreact to news)\citep{doh_lee_park_2024_rwp2406}. By contrast, higher-income (or more educated) respondents tend to have expectations closer to experts' and update somewhat more modestly.
Education level similarly correlates with expectation accuracy: those with a college degree generally forecast inflation or economic conditions more accurately and with less variance than those with only high school education. This could reflect higher financial literacy or numeracy, better access to quality information, or greater trust in expert sources. It may also reflect differential exposure to sources of misinformation or sensationalist news, which often target or resonate with certain demographics.
Political affiliation (though not an SES per se) has been shown to cause large expectation divergences: after recent U.S. elections, households aligned with the winning party become much more optimistic about the economy, while the other side becomes pessimistic, despite experiencing the same objective macroeconomy \citep{miansufi2023}. \citet{miansufi2023} document huge partisan gaps in expected economic growth after 2016, for example, driven by subjective narratives of the incoming administration's impact. Interestingly, they find these expectation differences had \textit{no corresponding effect on actual spending behavior on average}. This indicates that not all expectation biases translate into action, a point we will revisit. Nonetheless, it underlines how narratives (in this case political narratives) can sway beliefs differently for different groups.
Taken together, these facts motivate our focus on SES heterogeneity: it appears that lower-SES individuals both receive information differently and respond to it differently. One plausible mechanism is that \textbf{information sources} differ by SES. For example, surveys show that education level correlates with media consumption patterns; highly-educated people may rely more on newspapers or financial news, whereas less-educated individuals may rely more on television or social media for economic information. Podcasts as a medium span this divide to some extent (podcast listeners skew younger and more educated than the general population, but the content can vary widely). In any case, if polarized narratives are more prevalent in the media channels frequented by certain SES groups, those groups' expectations could systematically deviate due to narrative influence.

\subsection*{Narratives, Echo Chambers, and Inequality-Relevant Outcomes}
Why are we concerned with narratives and echo chambers specifically? Because there is evidence that media can causally affect beliefs and downstream outcomes. Classic studies by \citet{DellaVigna2007} showed that the introduction of the partisan Fox News channel in the late 1990s measurably shifted voting behavior: towns that gained access to Fox News saw a 0.4 to 0.7 percentage point increase in Republican vote share in 2000 \citep{DellaVigna2007}. This suggests persuasion effects of biased media are real. Fox News consumption has also been linked to altered beliefs about factual matters and policy. \citet{GentzkowShapiro2010} found that newspapers' slant is strongly tailored to local political leanings, implying consumers self-select into news that aligns with their priors \citep{DellaVigna2007}. Thus, media markets tend to create silos of information (a demand-driven echo chamber effect). Over time, this can reinforce and amplify belief differences among groups, including economic beliefs.
There is direct evidence that economic narratives can drive expectation errors. For example, \citet{ArmantierEtAl2016PriceIsRight} showed that exposure to more negative news about inflation (say, through certain news outlets) correlates with higher perceived inflation and higher inflation expectations among consumers, even controlling for actual price changes. Similarly, qualitative surveys (open-ended responses about economic concerns) indicate that people often echo narratives they hear in media (e.g., talk of recession or ''booming economy''), which then correlates with how they forecast their own prospects \citep{Shiller2017}. These findings hint that narrative content is a significant input into expectation formation.
Regarding \textbf{inequality-relevant margins}, the concern is that if one group persistently holds biased expectations, they may systematically underperform economically. For instance, consider a pessimistic narrative that convinces some households a recession is imminent. Those households might cut back spending or avoid investing in risky assets (like stocks or a new business) due to fear, whereas less influenced households might continue with normal or even optimistic behavior. If the recession narrative turns out false (no recession occurs), the pessimistic group has unnecessarily foregone opportunities (lower stock returns, etc.), potentially losing ground financially relative to the optimistic group. Over years, such effects compound. Another example: if some households overestimate future inflation (perhaps due to hearing hyper-inflationary rhetoric in their media bubble), they might demand higher nominal wages or shift savings to inflation-hedged assets (like real estate or gold). If inflation instead stays low and stable, those decisions could be suboptimal (e.g., not taking advantage of low inflation environment in fixed-income investments, or mispricing their wage demands leading to unemployment). In contrast, households with well-anchored expectations might make better financial plans.
While concrete evidence directly linking narrative-driven expectation bias to wealth outcomes is scarce (likely due to data limitations), this project intends to explore some \emph{proxy outcomes}. For example, we can look at whether regions with high echo chamber exposure saw divergences in household credit growth or default rates during certain periods, or whether individuals in those areas altered durable consumption in ways unrelated to fundamentals. If such patterns exist, it strengthens the case that narratives not only change beliefs but have real economic consequences.
Finally, the rise of podcasting as a medium for detailed discussions (including economic and political talk) provides a new window into narratives. Unlike traditional news, podcasts often involve long-form commentary, which may have stronger persuasive power through storytelling and repetition. The podcast landscape is also highly fragmented and polarized: some of the most popular news podcasts are overtly partisan or have strong ideological slants (e.g., those hosted by political figures or commentators), and their audience has grown rapidly in recent years[12]. According to Edison Research, monthly podcast listenership in the U.S. grew by 20\% from 2022 to 2024, with 48\% of those aged 12+ listening to a podcast in the past month[12]. As podcast consumption becomes mainstream, understanding its influence on economic perceptions becomes increasingly relevant.
In summary, the evidence suggests (1) expectation formation is prone to errors and biases; (2) these biases can differ systematically by group; (3) biased information sources (echo chambers) are a plausible cause of some biases; and (4) such biases could influence decisions with potential inequality implications. These points motivate our research design to isolate the causal impact of narrative echo chamber exposure on expectations and behavior.

\section{Related Literature and Gap}
Our project builds on and connects several strands of literature in macroeconomics, finance, and media economics:

\paragraph{Expectation Formation Biases:} The work of \citet{Coibion2015} (CG) introduced regression tests for information rigidities, showing that forecast errors are predictable by past forecast revisions in many contexts, contrary to full-information rational expectations. They interpret a positive error--revision correlation as evidence of \emph{sticky information} or under-adjustment. This provided a set of \textit{stylized facts} to calibrate models of imperfect information[2]. Subsequent research by \citet{Bordalo2018Diagnostic, Bordalo2020NE} (BGLS and coauthors) provided a behavioral alternative: \textit{diagnostic expectations}, wherein agents over-weight new information. This generates negative error--revision correlations (overreaction) at the individual level[1], which they document in survey data for various macro variables. Our proposal is directly inspired by BGLS's finding of overreaction: we seek to uncover one source of the ``diagnostic signals'' that might drive such overreaction in the wild (namely, narrative echo chambers).
Another perspective comes from \citet{Sims2003} and related rational inattention models, which assume agents have limited \textit{information-processing capacity}. Sims argued that it is plausible to model agents as optimally choosing what signals to pay attention to given a Shannon channel capacity constraint[5]. These models can generate inertia and sluggish updating as agents ignore small shocks. However, they typically would not generate \emph{systematic biases} (agents still use information efficiently given constraints). Our interest is in biases arising from \emph{the nature of information agents receive}. In that sense, our work is closer to models of \emph{noisy information} and even \emph{strategic communication}. Indeed, recent work by \citet{Andre2023biasedforecasts} finds that professional forecasters' published forecasts may be strategically biased (e.g., to gain attention or reputation), which can also produce patterns of over- or under-reaction[13][14]. While their focus is on deliberate bias by forecasters (supply side of expectations), our focus is on bias in the information environment of households (demand side).
\paragraph{Media Influence and Polarization:} A rich literature in political economy and media economics examines how media content is determined and how it influences beliefs. \citet{GentzkowShapiro2010} famously constructed an index of newspaper slant based on the similarity of phraseology to Republican or Democrat politicians. They found that newspapers' slant correlates strongly with the ideological bent of their local readership, consistent with outlets slanting news to cater to consumer preferences (``demand-driven slant'')[11]. Their model implied that if a newspaper deviated from the profit-maximizing slant even slightly, it would lose readers[15], underscoring how consumer demand can create self-reinforcing echo chambers. This is relevant for us because it suggests that if certain narratives (say doom-and-gloom about the economy, or rosy optimism) are in demand by particular demographic groups, the media ecosystem will supply those narratives, resulting in segmented pools of beliefs. We will attempt to measure something analogous for podcasts: e.g., is there segmentation such that certain regions or demographics predominantly listen to one type of narrative? Preliminary anecdotal evidence says yes—podcast rankings often show clusters of partisan or themed shows with devoted followings.
The causal effect of media on beliefs/behavior has been demonstrated in several contexts. Besides the Fox News effect on voting[10], \citet{GerberKarlanBergan2009} ran a field experiment providing people with free subscriptions to either the Washington Post or Washington Times and found effects on political knowledge and opinions. \citet{Baron2006} and others have theorized about supply-driven bias (e.g., owner's ideology influencing content). While our context is economic narratives rather than political news per se, these insights carry over. In fact, economic news is often intertwined with politics (consider debates on inflation causes, or unemployment and elections). \citet{ShapiroGentzkow2009} find that media biases can even affect expectations about the economy, not just political opinions. We also note the work of \citet{DellaVignaKaplan2007}, which specifically showed that exposure to a biased news source (Fox News) measurably shifted voting and possibly attitudes, implying that at least some viewers were persuaded rather than discounting the bias[10]. This persuasion likely operates through narratives and repetition.
In the digital era, algorithms and social media have intensified concerns about echo chambers and ``filter bubbles'' (see e.g., \citealp{Sunstein2001}). Our project is complementary, focusing on a somewhat less-studied medium (podcasts) that nevertheless shares features with social media: easy global distribution, niche targeting of content, and often unregulated or unedited monologues. According to a Pew Research report, about one in ten U.S. adults often get news from podcasts, and Republican-leaning listeners are more likely than Democrats to trust news from podcasts over other sources[16]. This suggests that podcasts might especially reinforce the beliefs of certain groups (possibly due to trust in hosts or a feeling of community).
We also draw on literature about narrative economics. \citet{Shiller2017narrative} argues that popular narratives (like ``Housing prices never fall'' in the mid-2000s) can drive economic booms and busts. These narratives often spread via media and word-of-mouth. Our empirical strategy can be viewed as an attempt to capture the spread of a narrative (via podcast popularity and content) and trace its impact on expectations. Recent stats from Edison Research confirm the growing reach of the podcast medium for disseminating narratives[12].
Furthermore, an emerging literature in behavioral macro uses \emph{textual analysis} to quantify narratives. \citet{Shapiro2020news} analyze news articles to create measures of sentiment or uncertainty that correlate with economic outcomes. We will similarly analyze podcast transcripts for narrative sentiment (e.g., measuring how positively or negatively a show talks about the economy, or whether it emphasizes certain themes like debt crisis'' orjob creation''). This overlaps methodologically with the literature on news sentiment indices (e.g., Economic Policy Uncertainty index by \citealt{BakerBloomDavis2016}) but applied to spoken word content. The gap is that few if any studies have used \textbf{podcast transcripts} and linked them to survey expectations data.

\paragraph{Household Expectations and Macro-Finance:} Our work relates to the growing macro-finance literature that uses household expectation surveys to understand economic outcomes. Studies like \citet{Bachmann2015} and \citet{D'Acunto2021} examine how household inflation expectations influence spending on durables and other decisions. Generally, if households expect higher inflation, some may spend more (to buy before prices rise), while others might hold off (if they think their real income will erode). The net effect is empirical. We will use such insights when looking at outcomes: for example, we might test if areas where households became unusually pessimistic (due to a narrative shock) saw declines in spending intentions as measured by survey or even retail sales data.
\citet{MianSufiKhoshkhou2022} (MSK) specifically study \emph{partisan bias in economic expectations}. They find that partisan divergences in expectations have grown over time (especially post-2008 and post-2016 elections) and are large—e.g., after 2016, Republicans expected far higher income growth than Democrats did[8]. However, they also find that these biased expectations did not lead to differences in actual spending between the groups[9]. This is both a cautionary note and a puzzle: it could be that some expectation responses are more like "cheap talk" or expressive (not truly believed, or not acted upon). Alternatively, constraints and habit might prevent rapid spending adjustment. For our study, we take this into account by not assuming a one-to-one mapping from beliefs to actions. We will interpret our results carefully: evidence of belief distortion is our primary focus, and evidence of action changes will be explored but may be subtle. The MSK result underscores the importance of looking at actual behavior data (like administrative spending or credit data) to verify whether belief changes matter economically.

\paragraph{Shift--Share Causal Inference:} Methodologically, our use of a shift--share (Bartik) design is informed by recent econometric advances. \citet{Borusyak2022shiftshare} formalize the conditions under which shift--share instruments yield causal identification. In particular, they clarify that one needs either \emph{exogenous shocks} (the "shift" terms are as-good-as-random and independent of the region-level exposure shares) or \emph{exogenous shares} (the exposure shares are randomly assigned, conditional on some controls)[17]. In our context, the shocks will be changes in podcast content popularity or slant over time, and the shares will be baseline audience measures by region. We'll argue for the plausibility of shock exogeneity: e.g., the content produced by podcasters on national issues is plausibly not influenced by local economic conditions in each region (especially if we focus on national news podcasts). We will also control for regional fixed effects and time trends to absorb obvious confounders. The literature also offers guidance on statistical inference for Bartik designs (e.g., \citealp{Adao2019}). We will likely implement the recent \textit{shock-level clustering} or the leave-one-out approaches recommended to get correct standard errors when a few large shocks or shares dominate.
Additionally, \citet{BorusyakHullJaravel2023JEP} provide a practical guide to shift--share instruments, emphasizing checks like ensuring no single share or single shock overwhelmingly drives results (which could violate the ``quasi-experimental'' assumption)[17]. We will follow these best practices (for example, check that our results are not driven solely by one or two very popular podcasts). If needed, we might use a procedure to orthogonalize shocks or exclude influential units to test robustness.

\paragraph{Other Related Work:} A few other references worth noting: \citet{DiFrancesco2024ABM} (Di Francesco and Hommes, 2024) develop an agent-based model in which some agents receive distorted (mis)information signals about financial asset payoffs and others learn from neighbors. They show that misinformation can spread in a network and delay price discovery, making markets less efficient[18]. This is conceptually similar to our idea that echo chambers (a network phenomenon) delay or distort the incorporation of true fundamentals into people's expectations, causing protracted belief errors. Our project can be seen as attempting to test, in survey data, whether such diffusion of distorted signals occurs in reality. Finally, our work is complementary to the burgeoning literature on measuring \textit{expectations dispersion} and its causes (e.g., \citealp{MankiwReis2002, Coibion2012}). By introducing narrative fragmentation as a potential cause of dispersion, we hope to add a novel angle to that literature.
\paragraph{Gap and Contribution:} In summary, the gap we aim to fill is at the intersection of these literatures: \textbf{no prior study has causally linked a quantified measure of narrative echo chamber exposure to a quantified measure of diagnostic expectation distortion (overreaction) in a micro-data setting, accounting for heterogeneous effects by SES.} We bring together techniques from media analysis (text mining of transcripts, audience measurement) with cutting-edge identification in macro (shift--share instruments on survey microdata). By doing so, we hope to demonstrate a concrete channel through which \emph{narrative-driven correlated noise} enters expectations. The novelty lies especially in directly measuring narrative exposure intensity ($Z^{\text{int}}$), narrative direction ($Z^{\text{dir}}$), and narrative fragmentation ($Z^{\text{frag}}$) and linking those to deviations from rational updating (the $\theta$ diagnostic overreaction parameter).
Crucially, we investigate the distributional dimension: earlier work either looked at professional forecasters or average households, but we will see whether the ``echo chamber effect'' is more pronounced among certain demographics. This has important policy implications—if, say, less-educated households are more prone to being misled by economic narratives, then targeted communication by central banks or authorities might be needed to counteract popular misperceptions in those groups. It also speaks to inequality: understanding whether information ecosystems can exacerbate economic inequalities by skewing expectations and decisions in certain populations.
Thus, our contributions will be: (1) providing evidence of narrative-induced expectation distortions (or showing absence thereof, which is equally informative); (2) introducing new measures of media narrative exposure for macro researchers; (3) highlighting SES-based information frictions as a component of heterogeneity in macro models; and (4) evaluating the macroeconomic significance of echo chambers (do they just create noise, or do they amplify shocks and affect outcomes?).
Next, we outline our conceptual framework and hypotheses in more detail, then our empirical strategy for identification and estimation of these effects.
\section{Conceptual Framework}
We conceptualize the relationship between narrative echo chambers and economic expectations as a causal chain of information flow and decision-making. Figure~\ref{fig:dag} presents a stylized directed acyclic graph (DAG) summarizing this framework. The key elements are: the supply of narratives in media, the local exposure to those narratives, the process of belief updating by households (with potential distortion parameterized by $\theta$), and the subsequent impact on household decisions and ultimately economic outcomes.

% TODO: CHECK THIS
% \begin{figure}[h!]\centering
%     \begin{tikzpicture}[->,>=stealth, node distance=1.5cm, thick]
%         \node[draw, rectangle, align=center] (narr) {Polarized Narrative\\Content (Podcasts)};
%         \node[draw, rectangle, below of=narr, yshift=-1cm, align=center] (exposure) {Local Narrative\\Exposure $Z_{r,t}$};
%         \node[draw, rectangle, below of=exposure, yshift=-1cm, align=center] (belief) {Belief Updating\\Distortion $\theta_{i}$};
%         \node[draw, rectangle, below of=belief, yshift=-1cm, align=center] (expect) {Expectations of\\Household $i$};
%         \node[draw, rectangle, below of=expect, yshift=-1cm, align=center] (decisions) {Household Decisions\\(spending, saving, etc.)};
%         \node[draw, rectangle, below of=decisions, yshift=-1cm, align=center] (outcome) {Inequality-Relevant\\Outcomes};
%         \draw[->] (narr) -- node[right]{e.g. partisan content} (exposure);
%         \draw[->] (exposure) -- node[right]{echo chamber intensity}(belief);
%         \draw[->] (belief) -- node[right]{over-/under-weight info}(expect);
%         \draw[->] (expect) -- node[right]{perceived constraints/opportunities}(decisions);
%         \draw[->] (decisions) -- node[right]{aggregate distributional effects}(outcome); % Moderators
%         \node[draw, ellipse, left of=belief, xshift=-3.5cm, yshift=-0.25cm, align=center] (ses) {Household SES\\(income, education)};
%         \draw[dashed,->] (ses) -- node[above]{moderates}(belief);
%         \draw[dashed,->] (ses) |- (decisions);
%         \node[draw, ellipse, right of=exposure, xshift=3.5cm, yshift=-0.5cm, align=center] (frag) {Media\\Fragmentation};
%         \draw[dashed,->] (frag) -- node[above]{shapes}(exposure);
%         \draw[dashed,->] (frag) |- (belief);
%     \end{tikzpicture}
%     \caption{Causal diagram illustrating how exposure to narrative echo chambers could distort belief updating and propagate to economic decisions and inequality. Dashed lines indicate moderating influences (e.g. socio-economic status and media fragmentation).}
%     \label{fig:dag}
% \end{figure}
Starting at the top: \textbf{Narrative content} is produced by various media outlets (podcasters, radio shows, news sites). In our context, consider the content of a set of popular economic or political podcasts. At any given time $t$, each show $j$ might be emphasizing certain themes or narratives. Some might be shouting recession alarms, others may be painting a rosy picture of economic growth, often colored by political ideology or other biases. We can think of each show's content as providing a signal about the state of the economy, but one that may contain \textit{bias or noise}. For example, a highly polarized show might selectively present facts to fit a narrative (“the economy is crashing due to policy X” or “the economy is booming thanks to policy Y”). In expectation-formation terms, these are \textbf{noisy signals} $s_{j,t}$ that agents observe.
Next, \textbf{local exposure $Z_{r,t}$}: not everyone listens to the same shows. We index regions by $r$ (e.g. Designated Market Areas or commuting zones) and time $t$. $Z_{r,t}$ represents some aggregate measure of how much the population in region $r$ is exposed to biased narratives at time $t$. We will operationalize $Z$ in multiple ways (intensity, direction, fragmentation as noted), but conceptually it captures the idea that in some regions/times, echo chamber narratives are more pervasive. This could be due to higher listenership of partisan podcasts, or a few especially influential local media voices, etc. $Z$ thus links the supply of narratives to the information set of households.
Crucially, we think of $Z$ as affecting the \textbf{signal that households receive}. In a rational expectations world, households would get unbiased signals of fundamentals (plus perhaps some random noise). In our world, households in a high-$Z$ environment get a signal that is biased or excessively noisy in a particular direction (correlated noise). We formalize the belief updating for an individual $i$ in region $r$ as follows: let $y^*$ be the true fundamental (e.g. next quarter GDP growth) that agents try to forecast. A household $i$ starts with prior belief $E_{i}(y^*| \text{prior})$. Then they receive new information. Under an echo chamber, the new information could be represented as:
% TODO: CHECK THIS
\[
    s_{i,t} = y^* + \underbrace{\varepsilon_{t}}_{\text{fundamental noise}} + \underbrace{\delta \cdot \mathbf{1}_{\{ i \text{ in narrative } X \}}}_{\text{narrative-specific bias}},
\]
where $\varepsilon_t$ is some normal noise of unbiased news, and the $\delta$ term is a bias that affects all individuals tuning into a particular narrative $X$. For example, if narrative $X$ is ``the economy is doing much worse than official figures show'', then $\delta$ might be a negative bias that subtracts a certain amount from the signal for all those listeners.
Household $i$ updates its expectation of $y^*$ upon seeing $s_{i,t}$. A rational Bayesian (with known signal distributions) would compute a weighted average of prior and new signal. A \textbf{diagnostic updater} as in BGMS, however, uses a distorted weighting: effectively $E_{i}(y^*|s) = E_{i}(y^*|\text{prior}) + \theta \big( s_{i,t} - E_{i}(y^*|\text{prior}) \big)$, where $\theta > 1$ means overreaction (overweighting the new signal relative to Bayesian weight) and $\theta < 1$ means underreaction (underweighting). Here we explicitly include $\theta$ as an individual-specific distortion parameter. In general, $\theta$ could be a function of the signal itself (Bordalo et al model $\theta$ increasing with the surprise magnitude), but as a first approximation we consider it a personal trait or group trait.
Our hypothesis (H1) is that $Z_{r,t}$, the narrative exposure, increases $\theta_{i}$ for individuals in that region/time. Intuitively, being in an echo chamber might hype up the perceived importance of signals (making them more ``representative'' in the BGMS sense). For instance, constant repetition of recession talk makes any negative data seem extremely representative of a coming downturn, so people overweight that data. Formally, one could derive from a psychological model that hearing confirming narratives increases the agent's confidence in the signal (or reduces perceived variance of noise), which leads them to put a higher weight on it. Thus, $Z$ could modulate the weight $\theta$.
Alternatively, $Z$ might directly tilt the signal ($\delta$ term above). In that case, even a rational updater would have a biased posterior because the input is biased. However, rational agents might discount known biases. If households are unaware of the bias or trust the source fully, they effectively treat $s_{i,t}$ as unbiased, yielding a biased expectation. This again can be captured by a $\theta$ interpretation: they effectively have $\theta=1$ (normal weight) but on a biased signal, which results in a predictable error.
In either story, the outcome is that expectations $E_{i,t}[y^*]$ differ systematically for high $Z$ vs low $Z$ environments, even conditioning on the true fundamental and prior. Echo chambers cause \textbf{distorted expectations}. Specifically, if a narrative is overly optimistic, exposed individuals will have overly optimistic expectations ($E - y^*$ positive); if overly pessimistic, expectations undershoot reality ($E - y^*$ negative). Our empirical focus is capturing these \textit{belief distortions} in the data via forecast errors and revisions.
Now consider \textbf{SES heterogeneity}. We include in Fig.~\ref{fig:dag} an SES node pointing (dashed lines) to both the belief updating and decision nodes. SES (say low vs high) can moderate the effect at multiple stages: \begin{itemize} \item \emph{Exposure stage:} SES might determine which narrative an agent is exposed to. For example, perhaps lower-SES individuals are more likely to listen to sensationalist economic podcasts or talk radio, whereas higher-SES might favor more staid financial news. If so, the content bias $\delta$ is SES-dependent. Also, lower-SES might have less access to diverse information, meaning their $Z$ is effectively higher (they're more ``all-in'' on one narrative). Higher-SES might consume multiple sources (lower effective $Z$ from any single narrative). \item \emph{Updating stage:} Even given the same exposure, SES might alter $\theta$. If higher education imparts more skepticism or statistical reasoning, high-SES agents might inherently have $\theta$ closer to 1 (they don't overreact as much). Lower-SES might trust the narrative more or have fewer alternative reference points, yielding higher $\theta$. This is what H2 posits. \item \emph{Decision stage:} SES differences in resources can buffer or amplify behavior changes. For instance, a wealthy person who believes a recession is coming might reduce spending slightly, whereas a poorer person might dramatically cut spending out of precaution (or vice versa: a rich pessimist might not need to change lifestyle, a poor pessimist might have to). Also, access to credit or assets differs: a low-income household that mistakenly believes inflation will spike might not be able to act on that belief effectively (they can't easily go invest in inflation hedges), whereas a high-income household could shift their portfolio. So even if beliefs shift equally, the realized actions could differ by SES. \end{itemize}
The \textbf{media fragmentation} moderator in Fig.~\ref{fig:dag} reflects how the structure of the media market in a region can affect exposure and belief variance. If a region has highly fragmented media (many small podcasts rather than one or two big ones), people may self-sort into different echo chambers. This would mean high overall $Z^{\text{int}}$ (lot of narrative listening) but split across conflicting narratives ($Z^{\text{dir}}$ might be both high positive for some and high negative for others). The net effect could be that the average expectation in that region is near rational (if biases cancel), but disagreement is high. By contrast, low fragmentation (one dominant narrative) might push the whole distribution of beliefs in one direction (high average bias but perhaps less disagreement internally). We will test this by examining dispersion of expectations as an outcome, as per H4.
Finally, \textbf{household decisions and inequality outcomes}: If expectations are distorted, households may make suboptimal decisions. These decisions aggregate up to things like consumption growth, saving rates, or default rates. Over time, if one group consistently makes worse decisions because of biased expectations, they might experience slower wealth accumulation or more financial stress, contributing to economic inequality. For instance, consider credit card debt: if a narrative leads some households to be overly optimistic about their future income, they might over-borrow now, leading to higher delinquency later predominantly in that group. Or housing investment: an echo chamber that downplays the risk of variable-rate mortgages might lead followers to take on such mortgages and suffer more when rates rise. These are speculative examples; our analysis will look at whatever outcome data we can link, but these mechanisms illustrate why an effect on inequality is plausible.
It is worth noting the possibility of \textbf{general equilibrium} feedback: if a large fraction of population holds distorted beliefs, their collective actions could actually influence the macro outcome $y^*$. For example, widespread pessimism can reduce consumption and investment, potentially causing a slowdown (a kind of self-fulfilling prophecy). Our study is mostly positive and partial equilibrium (we treat the macro fundamentals as given when measuring expectation errors), but one extension (discussed later) could explore whether narrative shocks have aggregate effects.
To test the conceptual links, we will estimate parameters that correspond to arrows in Fig.~\ref{fig:dag}: e.g., how much does higher $Z$ (\emph{exposure}) increase $\theta$ (\emph{belief distortion}) on average? How much more does a low-education individual react per unit $Z$ than a high-education one? Does high fragmentation increase the variance of $E_i(y^*)$ controlling for fundamentals?
In the next section, we turn to the empirical strategy detailing how we will measure these constructs ($Z$, $\theta$, etc.) and how we will achieve identification of the causal effect of exposure on expectations.
\section{Empirical Strategy: Identification and Estimation}
Our empirical strategy combines a \textbf{two-stage estimation approach} for identifying the effect of narrative exposure on expectation updating. In broad strokes:
\begin{enumerate}
    \item First, we measure households' \textbf{expectation updating behavior} – specifically, the degree of over- or under-reaction in their forecasts – using microdata on their expectations (from surveys). This gives us an estimate of the \emph{diagnostic distortion} parameter $\theta$ (or an analog, like the coefficient $\phi$ in a forecast error regression) for individuals or groups.
    \item Second, we relate this measured distortion to our \textbf{exogenous variation in narrative exposure} using a shift–share instrumental variable design. The shift–share instrument isolates variation in regional narrative exposure that is plausibly unrelated to local economic shocks, allowing a causal interpretation of its effect on expectations.
\end{enumerate}
We detail each part below, followed by the econometric specifications.


\subsection{Measuring Belief-Updating Distortions ($\theta$)}
To quantify the extent of diagnostic overreaction vs. underreaction, we will use the approach pioneered by \citet{CoibionGorodnichenko2015} and \citet{BordaloGennaioliMaShleifer2020}. The basic idea is to examine the relationship between agents' forecast revisions and their subsequent forecast errors. Define for individual $i$ (or a group) forecasting variable $y$ at horizon $h$ (e.g., inflation 1-year ahead):
\begin{align*}
    \text{Revision}_{i,t} & = F_{i,t}(y_{t+h}) - F_{i,t-1}(y_{t+h}), \\
    \text{Error}_{i,t+h}  & = y_{t+h} - F_{i,t}(y_{t+h}),
\end{align*}
where $F_{i,t}(y_{t+h})$ is the forecast made at time $t$ for outcome $y$ at time $t+h$, and $y_{t+h}$ is the realized outcome.

Under rational expectations with full information, $\mathbb{E}[\text{Error}_{i,t+h} \mid \text{Revision}_{i,t}] = 0$; past revisions should not predict errors (no systematic under- or overshoot). If we run a regression:
$$
    \text{Error}_{i,t+h}=\alpha+\phi\text{Revision}_{i,t}+\eta_{i,t+h},
$$
we expect $\phi = 0$ under the null of rational expectations. $\phi \neq 0$ indicates a bias:
\begin{itemize}
    \item $\phi > 0$: forecast errors tend to have the same sign as the revision. If $F$ was revised up, the error $y - F$ is likely positive (meaning $F$ was too low). This implies underreaction: they should have revised up more (they were still too low). This aligns with \textit{information rigidity/stickiness}\citet{GEMMI2026103868}.
    \item $\phi < 0$: errors are opposite to revisions. If forecast was revised up, error is likely negative ($F$ too high). This implies overreaction: they over-adjusted to the news. This is what diagnostic expectations predict\citet{GEMMI2026103868}.
\end{itemize}
In practice, we will estimate such regressions using our survey data. The NY Fed SCE provides panel data on expectations, so we can observe revisions at the individual level (e.g., person's inflation expectation in January vs June) and then see errors (using realized inflation). For some variables like personal income growth or spending expectations, realized outcomes are harder (we might use their next realization or ex-post self-reports). In other cases, we might use the cross-sectional mean as a proxy for realized or simply treat forecast of X and later forecast of X as iterative (this yields similar regression tests as CG 2015 approach).
One critical addition: we will allow $\phi$ to differ by exposure and by SES. Specifically, our goal is to estimate how $\phi$ (the degree of over/underreaction) varies as a function of narrative exposure $Z$ and SES. To do so, we can interact the revision with exposure measures. For example:
$$
    \text{Error}_{i,t+h}=\alpha+\phi_0\text{Revision}_{i,t}+\phi_1 (\text{Revision}_{i,t}\times Z_{r(i),t} )+\phi_2 (\text{Revision}_{i,t}\times 1\{\text{ low-SES} \})+\phi_3 (\text{Revision}_{i,t}\times Z_{r(i),t}\times 1\{\text{low-SES} \})+\text{FE}_{r,t}+\epsilon,
$$
where $r(i)$ denotes $i$'s region, and we might include region-by-time fixed effects $\text{FE}_{r,t}$ to absorb common shocks at that level (we discuss the fixed effect strategy later).
In this regression:
\begin{itemize}
    \item $\phi_0$ is the baseline reaction coefficient in low-exposure, high-SES scenario (if others are zero).
    \item $\phi_1$ captures how the reaction changes when exposure $Z$ is higher. We expect $\phi_1 < 0$ if echo chamber exposure induces more overreaction (i.e., pushes $\phi$ towards negative).
    \item $\phi_2$ would capture any general SES difference in reaction. We might expect $\phi_2 < 0$ if low-SES are more overreacting on average (or $\phi_2 > 0$ if low-SES underreact more).
    \item $\phi_3$ is the triple interaction: are low-SES in high-exposure areas especially prone to overreaction? This could be the case if both factors reinforce each other.
\end{itemize}
However, estimating such a granular interaction in one step might be ambitious given sample sizes. An alternative approach is two-step: first estimate individual-level (or region-level) $\phi$ coefficients, then regress those $\hat{\phi}$ on exposure and SES characteristics in a second stage. For example, one can estimate a separate Coibion-Gorodnichenko regression for each region (or each demographic group) and see how the coefficients correlate with average $Z$. But separate region estimates will be noisy unless sample per region is large.
Our solution is to incorporate the exposure measure directly in the forecasting equation through interactions, but carefully use instruments for $Z$ to ensure we isolate causal variation (since exposure $Z_{r,t}$ might be endogenous with forecast revisions or errors if, say, a bad local economy drives both narrative consumption and forecast errors). This leads to the shift–share IV.

\subsection{Exposure Indices $Z^{int}$, $Z^{dir}$, $Z^{frag}$}
We construct three main indices to quantify narrative echo chamber exposure for each region $r$ in each time period $t$ (likely monthly or quarterly):
\begin{enumerate}
    \item \textbf{Narrative Intensity ($Z^{\text{int}}{r,t}$):} This index measures how intensely the population in region $r$ is exposed to partisan or biased narratives at time $t$. Operationally, we define
          $$
              Z_{r,t}^\text{int} =\sum_{j \in \text{podcasts}}(\text{Share}_{r,j}^\text{base}  ) \times(\text{Intensity of narrative content in podcast } j \text{ at time } t).
          $$
          The term ${Share}_{r,j}^\text{base}$ is a predetermined weight reflecting the popularity of podcast $j$ in region $r$. We will derive these shares from historical data (e.g., relative download counts or Google search frequencies by region in a pre-study period). The idea is to fix how relatively exposed each region would be to each show's content if that show has a surge or drop.
          The ``Intensity of narrative content'' for podcast $j$ at time $t$ could be a measure like the fraction of that month's episodes that are about macroeconomic doom (or some measure of sensationalism). More simply, we might use something like the number of episodes or total listening minutes of podcast $j$ in period $t$ (nationally). The intuition: if a certain partisan podcast releases many episodes or gets a spike in listeners (perhaps due to some viral story), areas where that podcast has a big following (\textit{high share}) will experience a greater exposure shock.

          $Z^{\text{int}}$ therefore captures a shift--share: national shifts in content consumption, weighted by local predisposition to consume that content. A high $Z^{\text{int}}_{r,t}$ means region $r$ at time $t$ is predicted to have high narrative exposure (lots of biased content being consumed).

    \item \textbf{Narrative Direction ($Z^{\text{dir}}_{r,t}$):} This index captures the ideological or directional slant of the narratives predominantly heard in region $r$. We can define it similarly:
          \[
              Z^{\text{dir}}_{r,t} = \sum_{j} \text{Share}_{r,j}^{\text{base}} \times \Big( \text{Slant}_{j} \times \text{Intensity of $j$ at $t$} \Big).
          \]
          Here $\text{Slant}_{j}$ is a score for podcast $j$ indicating its political/economic bias (e.g., on a left-right spectrum or pessimistic-optimistic spectrum). We can derive $\text{Slant}_{j}$ from content analysis (for example, count phrases or sentiment that correlate with known partisan positions, or simply classify shows by known ideology of hosts).

          $Z^{\text{dir}}$ will tell us whether a region is more exposed to, say, a conservative economic narrative (positive value) or a liberal narrative (negative value) at time $t$. This is important because the nature of expectation bias might depend on direction. For instance, one narrative might consistently predict higher inflation (e.g., those worried about government deficits might always warn of inflation), while another might downplay inflation. If our outcome of interest is inflation expectations, $Z^{\text{dir}}$ could explain whether a region expects higher or lower inflation bias.

          We will likely interact $Z^{\text{dir}}$ with the sign of forecast errors to test if, for example, Republican-leaning narratives cause over-prediction of growth when a Democrat is in office (just as a hypothetical pattern).

    \item \textbf{Narrative Fragmentation ($Z^{\text{frag}}_{r}$):} This index is more static (or slow-moving). It measures how fragmented or concentrated the media landscape is in region $r$. One way to measure fragmentation is the effective number of popular podcasts. For example, compute a Herfindahl-Hirschman Index (HHI) of podcast audience shares in region $r$:
          \[
              \text{HHI}_r = \sum_{j} (\text{Share}_{r,j}^{\text{base}})^2.
          \]
          A high HHI means one or two shows dominate (low fragmentation), a low HHI means audiences are spread among many shows (high fragmentation).

          We can define $Z^{\text{frag}}_r = 1/\text{HHI}_r$ (effective number of sources) or something similar. Alternatively, $Z^{\text{frag}}$ could incorporate whether the top shows are ideologically diverse or all on one side. Perhaps measure the split between, say, right-leaning and left-leaning content consumed (50/50 would be fragmentation across ideologies; 90\% one side is low fragmentation in narrative viewpoint).

          This index does not vary (much) over time in our design, but it will be used in cross-sectional analysis and as an interaction term to test H4: does high fragmentation correlate with higher disagreement in expectations, controlling for other factors?
\end{enumerate}


We will standardize these indices for interpretability (e.g., in units of one standard deviation).
One might worry that these $Z$ measures could be endogenous: for instance, if a region experiences a local economic shock (say a factory closure), perhaps interest in certain economic podcasts spikes (people seek explanations), which would make $Z^{\text{int}}$ high \emph{and} could directly affect expectations (people become pessimistic due to the actual shock, not the narrative). This is why we employ a \textbf{shift–share instrument}. The instrument will rely on the idea that $\text{Share}_{r,j}^{\text{base}}$ (the predisposition of region $r$ for show $j$) is predetermined and the national variation in show $j$'s content or popularity is arguably exogenous to any one region's local economy. By interacting the two, we get predicted exposure that is not driven by local shocks.
Formally, let $N_{j,t}$ be a measure of narrative activity for podcast $j$ at time $t$ (e.g., number of episodes about economy, or a spike in listenership). Our \emph{first-stage} equation for exposure could be:
$$
    Z_{r,t}^\text{int} = \sum_j \text{Share}_{r,j}^\text{base} \, N_{j,t}
$$
We will instrument $Z^{\text{int}}{r,t}$ by its predicted value using national shocks $N$ as an orthogonal shock. In practice, one often cannot include all shocks separately due to limited $t$, so one uses the $Z$ itself as the instrument (since it's a linear combination of shocks). However, recent literature suggests doing shock-level clustering or using a leave-out mean for shares to avoid bias in coefficient estimates when using $Z$ as instrument and regressor. We will follow those recommendations (e.g., use $Z$ computed with leaving out region $r$ from share computation to ensure independence).
We will similarly construct instruments for $Z^{\text{dir}}$ (which is just $Z^{int}$ times slant sign of each shock) and possibly for interactions with SES if needed.
\subsection{Micro-level Estimating Equation}
Our main estimating equation at the individual level (suppressing some indices for clarity) will look something like:

\begin{equation}
    \label{eq:main}
    \text{Error}{i,t+h} = \beta_0 + \beta_1 Z^{\text{int}}{r(i),t} + \beta_2 \left(Z^{\text{int}}{r(i),t} \times \mathbf{1}{i\text{ low-SES}}\right) + \gamma \text{X}{i,t} + \mu_r + \tau_t + \varepsilon_{i,t+h}.
\end{equation}

Here, $\text{Error}_{i,t+h}$ might be signed such that positive means an over-prediction (forecast too high). The key coefficients:

\begin{itemize}
    \item $\beta_1$ captures how average forecast errors respond to echo chamber exposure. A significantly nonzero $\beta_1$ would indicate a bias. For example, $\beta_1 > 0$ might mean: high exposure leads to positive errors (overestimation of economy) on average, which could happen if the dominant narrative is over-optimism. $\beta_1 < 0$ could indicate high exposure leads to underestimation (perhaps a doom narrative).
    \item $\beta_2$ then indicates differential effect for low-SES. If $\beta_2$ is negative (and say $\beta_1$ is positive), that would mean low-SES have \emph{less} of the positive error or maybe negative error—i.e., they are more pessimistic for same exposure, or vice versa depending on signs.

\end{itemize}
However, the above equation might be too naive as it doesn't include the forecast revision. More appropriately, we might embed the exposure effect into the error–revision framework:

\begin{equation}
    \label{eq:cg_extended}
    \text{Error}{i,t+h} = \alpha + \phi ,\text{Revision}{i,t} + \psi Z^{\text{int}}{r,t} + \kappa \left(\text{Revision}{i,t} \times Z^{\text{int}}{r,t}\right) + \lambda{r,t} + \epsilon_{i,t}.
\end{equation}
In \eqref{eq:cg_extended}, $\lambda_{r,t}$ is a full set of region-by-time fixed effects. The inclusion of $\lambda_{r,t}$ means we are only exploiting within-region-at-time variation. How do we get such variation? Essentially through SES or other individual characteristics if they interact with exposure differently. Because if $Z_{r,t}$ is common to all individuals in region $r$ at time $t$, then $\lambda_{r,t}$ would soak up $\psi Z_{r,t}$ term fully. Thus, to identify $\kappa$ (the main coefficient of interest, effect of exposure on revision-error relationship), we rely on something like:
$$
    \text{Error}_(i,t+h)=\phi \text{Revision}_(i,t)+ \kappa (\text{Revision}_(i,t) \times Z_{r,t} \times 1{low-SES})+\lambda_{r,t}+\text{other interactions}+\epsilon.
$$
In words, with region-time FEs, we identify off differential slopes by SES within the same region and time. For example, we ask: in a given region-month, did low-education respondents' forecast errors vs revisions look systematically different from high-education respondents' in a way that correlates with the region's exposure level? This triple difference (low vs high education, high vs low exposure, before vs after revision) is complex but feasible. We will likely implement it as an interaction IV: use something like $(Z_{r,t} \times \mathbf{1}{\text{low-SES}})$ as an instrument for the interacted term in the CG regression.
Alternatively, we may aggregate to region-level or region-SES cell level. For instance, compute $\hat{\phi}{r, \text{low-SES}}$ and $\hat{\phi}{r, \text{high-SES}}$ from separate CG regressions by group, then regress those on $Z_r$ (with instrumentation). This two-step might be more transparent: “does high exposure correlate with more negative $\hat{\phi}$ (overreaction) for low-SES group?” etc. The challenge is standard errors and some errors-in-variables (we can possibly do a weighted IV on those estimates with known variances).

Another outcome we will examine is \textbf{dispersion of expectations}. For each region-month, we can compute the cross-sectional standard deviation of, say, inflation forecasts among respondents. We hypothesize that $\text{Var}(E_{i}[y])$ is higher when $Z^{\text{frag}}$ is higher, or when $Z^{\text{dir}}$ is polarized. We can test:
$$
    \text{Dispersion}_(r,t)=\rho_0+\rho_1 Z_r^\text{frag} + \rho_2 |Z_{r,t}^\text{dir}| + \text{controls}+u_{r,t}
$$
$|Z^{dir}|$ (absolute slant exposure) could lead to polarization between those who buy into it and those who don't, increasing dispersion. Or two offsetting narratives (one far left, one far right in area) also raise dispersion. We'll explore dispersion regressions as supplementary tests.

\subsection{Shift--Share Instrument and Identification}
The key identification challenge is that narrative exposure $Z_{r,t}$ might be correlated with unobserved factors that also affect expectations. For example, consider a scenario: a region is hit by a negative economic shock (factory closure). This might make people in that region both more pessimistic in expectations (for real reasons) and also more likely to consume negative narratives (e.g., they tune into podcasts that confirm their feeling the economy is bad, or local radio harps on the downturn). This would create a spurious relationship between $Z$ and expectation errors (they have pessimistic expectations which might look like overreaction to something, and $Z$ is high, but not because $Z$ caused it).
Our shift–share instrument addresses this by leveraging variation coming from the \textit{outside shifts} $N_{j,t}$ (national podcast trends) rather than local shocks. The identifying assumption is that the national fluctuations in podcast content/popularity are orthogonal to region-specific expectation shocks, conditional on fixed effects. For instance, say a prominent podcaster decides to do a series on "the coming recession" in mid-2024 due to some national news or personal choice. This will raise $Z^{int}_{r,t}$ more for regions where that podcaster has a big audience. As long as we assume the podcaster's decision wasn't caused by something specifically happening in region $r$ (plausible if it's a national show), then the differential exposure is as-good-as-random across regions in terms of local economic conditions.
We will include time fixed effects $\tau_t$ (or at least month-year dummies) to soak up any pure time trends (like overall podcast listening growth or nationwide sentiment changes). We include region fixed effects $\mu_r$ to handle any time-invariant differences (some regions always more optimistic, or always more podcasty, etc.). Then $Z$ variation left is essentially the interaction of those base shares with time-specific shocks. By instrumenting $Z$ with those interactions, we isolate the causal part.
We will test the validity of the instrument in several ways: \begin{itemize} \item Check balance: ensure that pre-treatment characteristics of regions (like average historical expectations or socio-demographics) are not correlated with the instrument (shock exposure). \item Placebo tests: apply the same IV strategy to a pre-sample period where we should see no effect (e.g., use base shares from pre-2018 with podcast shocks in 2019-2020 to predict expectation changes in 2015-2016—should be null because 2015-16 expectations couldn't be caused by later podcast content). \item Use the methodology of \citet{Adao2019} to test for correlation between shares and shocks. They suggest regressing shares on some function of shocks or vice versa to see if any pattern emerges. Ideally, each shock $N_{j,t}$ should not be disproportionately loading on shares from systematically different regions (that would be like one shock targeting specific region types). \item Over-identification: if we treat each major podcast shock as an instrument separately, we could test over-ID (though many shocks likely). \item Exogeneity to local economic outcomes: we can check if our instrument predicts local economic variables (like local unemployment changes) when it shouldn't. If we find that, say, a high exposure shock coincides with a local unemployment spike, that suggests instrument may be picking up that local shock effect (bad). \end{itemize}
In the second stage, we will be careful in interpreting coefficients. Due to the presence of fixed effects and interactions, the effective instrument is $\left(Z^{int}_{r,t} \times \mathbf{1}{i\text{ low-SES}}\right)$ for the triple diff perhaps. We will implement a suitable IV estimation (possibly a multi-dimensional one where we instrument each endogenous term with a corresponding shift–share).
Finally, we note an \textbf{intent-to-treat (ITT) vs treatment-on-treated (TOT)} issue. Our analysis at the region level treats the entire region's population as having exposure $Z_{r,t}$. But of course, not everyone in region $r$ listens to podcasts or to the particular narratives. Thus, our $\beta_1$ (for example) should be interpreted as an \emph{average effect on the population}, not on listeners specifically. It's an ITT effect of living in a high-exposure region. This likely understates the effect on actual treated individuals (podcast listeners). If only, say, 50\% of people are actually consuming the narrative, then a small shift in average expectations might hide a larger shift among that 50%. This attenuation is a classic issue when treatment is measured at aggregate level and not every individual complies with treatment.
We will discuss ways to address this: - We can attempt to measure the fraction of actual listeners in each region (e.g., from external data like Edison Research which might give fraction of population that listens to podcasts weekly). If we have that, we could scale our coefficients or at least say: if X\% are listeners and we find Y effect on average, then the effect on listeners is roughly Y/X. - In surveys, sometimes they ask how often you get news from source X. If SCE or other data has any question on media usage, we could use that to subset or to do an IV at individual level (like instrument individual's listening by region instrument). - As an extension, we might incorporate a model for partial compliance: treat $Z_{r,t}$ as instrument for an unobserved individual-level variable indicating actual narrative consumption. But without individual media data, that remains abstract.
Given these limitations, we will primarily interpret our findings as demonstrating a \textbf{population-level effect}: e.g., a one standard deviation increase in narrative exposure (due to an exogenous national narrative shock) causes, say, a 0.2 st.dev. shift in the average forecast error or a 0.1 increase in the overreaction coefficient $\phi$. We will note that the effect on actual listeners would be larger (if only a subset is affected, their shift must be more pronounced to move the average).

\subsection{Additional Tests and Variation}
Apart from the main causal tests, we will leverage a few supplementary analyses: \begin{itemize} \item \textbf{Directional effects:} Use $Z^{dir}$ to see if, for example, right-leaning narrative exposure predicts a bias in one direction of forecast error (maybe overestimating economic performance under a Democratic president, consistent with partisan pessimism, etc.). This can help confirm that it's indeed narrative content, not just any content, causing biases. \item \textbf{Time-series event studies:} Identify specific episodes where a narrative shock occurred (e.g., a popular podcast released a sensational series at a known date) and see if expectations in high-share regions moved differently around that time compared to low-share regions, controlling for national trend. This event study approach can visualize the timing alignment, strengthening causal claims. \item \textbf{Robustness to alternate $Z$ definitions:} e.g., define exposure by including talk radio or other media if data allow, or use different base period for shares to ensure results not sensitive. \item \textbf{Robustness to including local controls:} We will attempt to include controls for local economic conditions (like region unemployment rate, etc.). In a shift–share IV, including endogenous controls is tricky, but we can check OLS with controls to see if coefficient drops, as a gauge of omitted variables. \item \textbf{Dispersion regressions:} as mentioned, test H4 by regressing disagreement on fragmentation or on $Z$. \end{itemize}
Having outlined the strategy, we now describe the data that will enable this analysis, and the practical feasibility of obtaining and merging these data.
\section{Data}
Our analysis will combine data from several sources:
\subsection{Narrative Exposure Data (Podcasts and Media)}
To measure narrative echo chamber exposure, we require data on podcast content and consumption by region:
\begin{itemize}
    \item \textbf{Listen Notes API / Podcast Index:} We will use the Listen Notes API (a comprehensive podcast database) to gather metadata on a broad set of podcasts. This includes:
          \begin{itemize}
              \item Podcast show information: title, category/genre (e.g., News, Politics, Business), publisher, etc.
              \item Episode listings with timestamps and possibly brief descriptions.
              \item Importantly, Listen Notes sometimes provides popularity data or ``rank'' of podcasts, and it has a global search engine that can be queried for keywords.
          \end{itemize}
          Using this, we will identify a subset of podcasts that regularly cover macroeconomic or political-economic topics. For example, podcasts like \textit{Planet Money}, \textit{Marketplace}, \textit{Macro Musings}, or more partisan ones like \textit{Pod Save America} (liberal) or \textit{The Ben Shapiro Show} (conservative, which often discusses economy among politics). We will curate a list of perhaps the top 100-200 such podcasts by audience size.
    \item \textbf{Transcripts/Text Data:} For those podcasts, we will attempt to obtain transcripts. Many major podcasts provide transcripts on their websites (or we might use speech-to-text on episodes). The transcripts allow content analysis. We will parse transcripts for economic keywords (``inflation'', ``jobs'', ``recession'') and sentiment or tone. We can score each episode or each podcast-month by how pessimistic vs optimistic it is about the economy, or whether it blames certain actors, etc.
    \item \textbf{Slant and Bias Measures:} Leveraging either external classifications or the text analysis, we will assign each podcast a slant rating. One method: check if the podcast is referenced in known bias databases (some organizations rate media for bias). Another: use the transcript text and measure its similarity to Republican vs Democrat politician speech (like Gentzkow and Shapiro did for newspapers). Or simply classify them manually if obvious (e.g., host affiliation).
          For economic tone, we might use a dictionary approach: count negative words related to economy or look at sentiment specifically within economic context.
    \item \textbf{Regional Popularity (Shift–Share Shares):} We need $\text{Share}_{r,j}^{\text{base}}$, the baseline popularity of podcast $j$ in region $r$. Ideally, we would have data on podcast audience by region. This is tricky since such data is proprietary (held by Spotify, Apple, etc.). However, we can approximate using indirect measures:
          \begin{itemize}
              \item \textit{Google Trends:} We can use Google Trends data on podcast name searches by metro area. For example, how frequently people in different cities search for ``Ben Shapiro podcast'' or related terms. If a region has a higher relative search index for a podcast, that suggests a larger audience share.
              \item \textit{Twitter/Facebook:} If data permits, look at social media followers of local accounts or mentions by location.
              \item \textit{Regional media surveys:} Sometimes Pew or others report what % of each region listens to podcasts, though not show-specific.
              \item A simpler assumption: use demographic info. Many podcasts target known demographics (e.g., some shows have younger urban audiences, others older rural). If detailed data is unavailable, we might define share proxies based on demographic similarity. But this is less ideal.
          \end{itemize}
          We will likely proceed with Google Trends as a feasible route. For each of our candidate podcasts, we obtain a relative popularity index for, say, the top DMA (Designated Market Area) or state. This gives a matrix of popularity. We then normalize each region's shares to sum to 1 across podcasts (or at least relative).

          Note: Google Trends gives relative interest, not absolute. We might need to calibrate such that $\sum_j \text{Share}_{r,j} =$ fraction of people who listen to those podcasts. If 40\% of US adults listen to podcasts monthly (approx), and of those maybe 20\% are news/politics, we can scale accordingly. This isn't crucial for regression (scaling doesn't affect IV consistency), but for interpretation of how many are treated, it helps.

    \item \textbf{Time Variation in Content ($N_{j,t}$):} As the shift (shock) variables, we will quantify how each podcast's content changes over time. Options:
          \begin{itemize}
              \item Measure of output: e.g., number of episodes in month $t$ that discuss economy or length of those episodes.
              \item Measure of tone: average sentiment or bias of podcast $j$ in month $t$. If a normally mild show suddenly becomes very pessimistic in a particular month, that's a shock in narrative content.
              \item External events: identify big news events and see how different podcasts react differently. But for instrument, we prefer aggregate measures that we treat as exogenous.
              \item Simpler: use overall episode frequency as a proxy for podcast being ``active''. Or their position on the charts (if a show jumps in rank, implies more listeners nationally).
          \end{itemize}
          We will likely construct $N_{j,t}$ as: (a) the monthly download rank or trending score of podcast $j$, if available via API or web scraping podcast charts; or (b) failing that, the number of tweets referencing that podcast (some services track this) as a proxy for popularity.

          Another angle: use the transcripts to measure narrative shock. For example, a topic model might identify a narrative theme (like "imminent recession" theme) and we can measure how much each podcast uses that theme over time.
\end{itemize}
We acknowledge that \textbf{data quality and availability for podcasts can be a limitation}. Listen Notes claims to have a global podcast index with some popularity metrics (they have an internal ``Listen Score'' metric). We will validate our exposure measures against any known facts (like if our data says New York has more of some podcast than LA, does that align with intuitive expectations or any survey evidence?).
There is also Edison Research's \emph{Infinite Dial} report which sometimes breaks down podcast listening by region or city (if we can get that, it could help calibrate shares). Pew surveys could tell us the overall % of people listening in rural vs urban etc.
Despite potential noise in these measures, as long as the measurement error in shares is not systematically correlated with outcomes, our IV should still be okay (though less powerful). We will conduct robustness by trying alternative share constructions.
\textbf{Listen Notes Data Quality:} Listen Notes is comprehensive (over a million podcasts indexed), but its metadata might be incomplete for smaller shows. However, we are focusing on more popular ones which likely have good info. One concern: some podcast networks might not publish precise data. We might rely on relative measures like ranks or episode counts.
In summary, through a combination of Listen Notes data, Google Trends, and text analysis, we will assemble a panel dataset of narrative exposure indices $Z$ by region-month.
\subsection{Expectations and Household Microdata}
For expectations, our primary source will be the \textbf{New York Fed Survey of Consumer Expectations (SCE)} microdata. Key features: - It's a monthly rotating panel (each respondent stays for up to 12 months). - It is nationally representative (internet-based sample). - It asks quantitative expectations on several domains: inflation (1-year and 3-year ahead), income growth, spending growth, house price growth, and the perceived probability of things like unemployment rise, etc. It also asks about uncertainty and some other behaviors (e.g., probability of missing a debt payment, plans to buy a car or house, etc.). - It collects demographics: age, gender, education, income, region (they have region identifiers, possibly state or at least census division, maybe even zip code or county in restricted data). - We will need region identifiers to merge with $Z_{r,t}$. We hope to get at least the state or CBSA (core-based statistical area). If only broad region (Northeast, Midwest etc.), that's too coarse. Prior literature using SCE sometimes has state. - SCE sample size: about 1300 per month. At any time, maybe ~100 per census region. Not huge for state-level, but possibly workable at division level or for large states. - We could aggregate to larger regions (like DMA or state groups) if needed.
We will use SCE to measure individuals' forecasts and back out revisions and errors: - We have their forecast at time $t$ for, say, inflation over next 12 months. Then a year later we see actual inflation and compute error. - If a respondent doesn't stay a whole year, we can't get their error. But we can look at group average errors (i.e., did the average in region r at time t overshoot relative to actual? But actual is national inflation, which is common to all – so error differences across regions in inflation would solely come from forecast differences). - Actually, for inflation or national variables, the actual is same for everyone. So any difference in error by region is entirely due to differing forecasts. That is fine since forecast differences are what narrative would cause. (We just have to account for the fact that one common actual means any aggregate shock affects all, but differential exposure could cause differential forecast). - For income or household-specific outcomes, the "actual" might differ by individual (someone actually got a raise or not). SCE might ask them later what their income is. If not, we might use self-reported realized or assume actual = actual macro stat (not perfect). - Alternatively, we focus on inflation expectations as our main dependent variable because it's clearly defined and common. And maybe unemployment or growth expectations if available.
We also have the University of Michigan's Surveys of Consumers (MSC), which is older and has long time-series, but that data at micro level is harder to get and has small sample (500/month) and only at broad region maybe. So SCE is better due to richer individual panel.
Other optional data: - \textbf{Consumer Expenditure Survey (CES):} If we wanted to see actual spending behavior changes, CES could be used (it's a rotating panel of spending diaries). We could see if high-exposure households cut spending differently. But linking to narratives may be complex. And CES sample by region might not be large. - \textbf{Home Mortgage Disclosure Act (HMDA) data:} gives all mortgage originations. Perhaps use it to see if, in high narrative areas, fewer mortgages are taken when narrative is pessimistic, etc. But that could be many confounds. - \textbf{Nielsen scanner or retail sales data:} Possibly use something like spending by region (if we have something like credit card data by region and time, we could test if consumption responds to narrative shocks). - \textbf{Fed Consumer Credit Panel (Equifax data):} if accessible, could see if credit card balances or delinquencies change by region after narrative shocks.
These are all extensions and may be beyond scope if data access is limited. The core analysis will use SCE for expectations (outcome) and the podcast-based measures (treatment).
We will carefully address \textbf{data privacy and quality}: The SCE microdata might require special permission, but as a PhD student I can likely get access (NY Fed often collaborates with researchers or they might have a public microdata release with scrambled IDs). If not SCE, an alternative is the Michigan Survey microdata (available through ICPSR). Or possibly the Survey of Consumer Finances panel (though that's triennial, not high freq). SCE is best due to monthly frequency to align with media monthly.

\subsection{Summary of Merged Dataset Structure}
The final panel for analysis might look like:
\begin{itemize}
    \item Unit of observation: individual $i$ in month $t$.
    \item Key individual vars: forecast of inflation (etc.), revision since last, demographic dummies, possibly their own media usage if available.
    \item Merged region vars: $Z^{int}{r(i),t}$, $Z^{dir}{r(i),t}$, $Z^{frag}_{r(i)}$ (these being IVs or treatments as appropriate).
    \item Fixed effects: time (month-year), possibly region or region-time if differencing.
\end{itemize}

We might also collapse some analysis to region-month level:
\begin{itemize}
    \item Then observation is region-month, with outcomes like average forecast, average revision, dispersion, etc., and treatments $Z$.
    \item That reduces N to (regions * time periods). If region = maybe 50 states, time ~ 60 months (5 years of data), that's 3000 obs, decent. If region = 9 census divisions, 60 months, 540 obs (maybe too small). We'll likely define region as around 50 states or 30 largest metro areas (since we have city-level shares possibly).
    \item If using state level, SCE might not have enough per state to get stable mean. We might cluster small states together (like DMAs or divisions).
\end{itemize}

Given these considerations, an intermediate approach: define perhaps 20 ``media regions'' (like grouping states by media markets) to ensure each has ~50 respondents per month to smooth means. The exact definition will depend on how fine the Google Trends can measure shares.

\subsection{Data Limitations:}
We should be transparent that measuring $Z$ is novel and thus subject to error. We might conduct a survey ourselves (small scale) to validate: ask people which podcasts they listen to and correlate with our predicted region exposure. If feasible, this would bolster credibility.
Also, the SCE sample though decent, may limit power in detecting some effects (especially triple interactions). We will use all tricks to increase precision: panel structure (difference out individual fixed effects if possible for some outcomes), and focusing on the strongest signals (maybe large narrative shocks events).
The Listen Notes data is quite large but manageable; we will likely preprocess transcripts to extract features rather than keep full text for all episodes (which could be many thousands of hours). We'll also focus on the top N shows to ease that.
Finally, merging by region requires careful alignment (e.g., Google Trends might be by DMA, we have to map SCE respondents to DMA via their zip code or similar). If SCE doesn't have zip, perhaps only state, we map as best as possible.
In summary, while data assembly is a heavy task, it is doable with modern text analysis and data from public APIs. The outcome is a unique dataset enabling our analysis.

\section{Threats to Identification and Robustness Checks}
We foresee several potential threats to a causal interpretation of our results and outline how to address them:

\paragraph{Endogenous Media Consumption (Demand for Slant)}
It's possible that people choose their media based on their pre-existing beliefs or local economic conditions. For example, pessimistic individuals might seek out pessimistic podcasts. This reverse causality would bias our results.

\textit{Mitigation:} Our shift–share IV is designed to account for this by using predetermined audience shares and national-level shocks that are not driven by individuals' current beliefs. We will also control for baseline differences in beliefs: e.g., include prior expectations or long-run expectations as controls to ensure we're capturing changes, not levels. Additionally, we can test for selection by checking if, before a narrative shock, high-share and low-share regions were on parallel trends in their expectations (placebo test).

Confounding by Unobserved Regional Shocks: Regions might experience unique shocks (economic or otherwise) that coincide with narrative exposure changes. For instance, a hurricane hitting region $r$ could both depress local outlook and also cause more media consumption (people stay home listening to news).

\textit{Mitigation:} We include region-by-time fixed effects in some specifications to flexibly absorb any common factor at that level. In the IV, the identifying variation is differential across SES within region-time, which should net out any region-wide shock. We will also incorporate controls for observable local economic indicators (unemployment rate, etc.) as a check; if adding those hardly changes results, it's reassuring. Moreover, using multiple instruments (different shows' exposure) can help—if an unobserved shock were the cause, it likely wouldn't align perfectly with the specific pattern of which podcast's shock affected which region.
National Trends and Omitted Global Shocks: A big national event (e.g., a federal policy change or a pandemic onset) can shift everyone's expectations and also become a hot topic in all media. This could create a spurious correlation between media content and expectations (both driven by the event). \textit{Mitigation:} Time fixed effects $\tau_t$ will capture any uniform national shift in expectations. Our identification relies on differential impact of that event via media. For example, during COVID-19, all regions got hit (time FE captures general pessimism), but perhaps some narratives (say conspiracy podcasts) pushed extreme expectations in some areas more than others—that differential is our focus. We will explicitly analyze such events: if our instrument is just picking up the direct effect of COVID on certain areas (maybe areas with older population got more pessimistic independent of media), that would violate exogeneity. We will check if excluding major crisis periods (like Mar-Apr 2020) affects results. Also, we might directly control for event intensity (COVID case rate by region, etc.) to ensure media effect isn't proxying for differing real impact of the event.
Ecological Fallacy (Aggregate measure, individual outcome): We assign an exposure $Z_{r,t}$ to all individuals in a region, but not everyone is actually treated. This can cause two issues: attenuation (as discussed) and potential misattribution if the composition of who listens differs from who responds in survey. \textit{Mitigation:} We interpret coefficients as ITT. We will attempt to gauge the fraction treated (podcast listeners) and discuss how to scale effects. To ensure we're not just capturing that our respondents aren't even listeners (so why would their expectations move?), we rely on social spillovers or second-hand effects—those exist (people talk to others who do listen, etc.), but weaker. We will try a subgroup analysis: e.g., younger respondents (more likely podcast listeners) should show stronger effects than older ones if mechanism is direct listening. If we see that pattern, it supports the interpretation. If not, it might suggest a broader community-level effect (which is still interesting).
Measurement Error in $Z$: Our constructed exposure indices are noisy proxies of true narrative exposure. This classical measurement error would bias coefficients toward zero (attenuation), making it harder to find effects. \textit{Mitigation:} The IV actually helps here too, since it's based on the same data but focusing on variation with better signal-to-noise (the aggregated shock). We will improve measurement by smoothing (e.g., using a 3-month moving average of $Z$ if needed) or dimension reduction (maybe use principal component if multiple correlated narrative indices). We'll also do robustness: e.g., define $Z$ using a different base period or excluding some podcasts and see if results hold, to ensure it's not driven by one outlier measurement.
Exogeneity of Shift–Share Instrument: For the Bartik instrument to be valid, two conditions: \textit{exogeneity of shocks} (the $N_{j,t}$ are unrelated to target outcome except through shares) and/or \textit{exogeneity of shares} (the $\text{Share}_{r,j}^{base}$ are as-good-as-random across $r$ after controlling for fixed effects). In our context, one worry is that the distribution of podcast popularity (shares) is not random—e.g., maybe conservative podcasts are popular in areas that also historically have certain economic volatility. Or shock exogeneity: maybe a particular podcast's content surge is itself a reaction to a localized phenomenon (though less likely as they target national audiences). \textit{Mitigation:} We will apply tests as recommended by \citet{Borusyak2022shiftshare}. For example, we can test for correlation between the share-weighted sums of pre-period characteristics and the shock series. We might also use the procedure of creating pseudo-shocks to test if any structure in shares would create bias. If a particular podcast (shock) is dominant and potentially problematic, we can exclude it and re-run (jackknife exclusion). If results remain, that increases confidence. Furthermore, if many shocks are small, the law of large numbers helps validity\citep{borusyakhull2025}; if one shock is large (like one show contributes a lot), we focus on showing that show's pattern still fits exogeneity (maybe through narrative that it wasn't targeting specific regions).
Local Spillovers and SUTVA violations: If one region's narrative exposure can spill over to another (e.g., people in neighboring region also listen), or if people travel, it could blur treatment. We assume region $r$ exposure mainly affects region $r$. With online media, geography is less tight, but our region definition might approximate communities of interest. \textit{Mitigation:} Use broad enough regions that cross-listening is minimal (if using state, probably okay; if using city, someone might commute from suburb which is another region in data). Also, if exposure is national but we attribute regionally, that's okay for instrument since national part is in time FE, only differential part matters. We will be cautious interpreting spatial differences.
Multiple Hypothesis / p-hacking concerns: We have several outcomes (inflation expectations, income expectations, etc.) and multiple indices ($Z^{int}, Z^{dir}$). There is a risk of cherry-picking significant results. \textit{Mitigation:} We pre-specify inflation expectation as primary outcome (since it's well-measured and widely relevant). Other outcomes will be secondary. We will apply a significance threshold and possibly adjust for multiple tests if needed (e.g., if we test 5 hypotheses, ensure the key ones meet appropriate significance or report q-values). But given this is exploratory to a degree, we'll focus on consistent patterns across measures rather than individual p<0.05 in isolation.
Each of these potential issues will be discussed in the final paper, with results of robustness checks (e.g., a table showing that controlling for local unemployment or using alternative instruments yields similar coefficients, bullet-proofing our interpretation).
\section{Expected Contributions and Extensions}
If successful, this project will make several contributions:
\begin{itemize}
    \item \textbf{New Evidence on Narrative-Driven Expectations:} We will provide, to our knowledge, the first causal evidence that exposure to polarized narratives (as proxied by popular podcasts) can distort household macroeconomic expectations. This moves beyond correlation to identify a concrete mechanism for expectation formation biases. It enriches behavioral macro by demonstrating that not only do people deviate from rational expectations, but specific identifiable information channels contribute to those deviations. In doing so, we empirically validate aspects of the \textit{diagnostic expectations} framework in a real-world setting: showing that narrative-induced overreaction is observable outside of lab or professional forecaster data.
    \item \textbf{Heterogeneity in Diagnostic Expectations:} We shed light on how diagnostic biases might differ across demographic groups. Prior literature often models a representative agent with a bias, or compares professionals vs households. Our findings will highlight SES-based heterogeneity: for example, we might find that less-educated households exhibit a larger diagnostic $\theta$ (more overreaction) when subjected to echo chambers, whereas more-educated are closer to rational (or even underreact). Such insight can inform heterogeneous agent models of expectations, suggesting that one size (bias) does not fit all. It can also inform policymakers that communications might need tailoring: e.g., central banks might have to counteract overly pessimistic narratives prevalent among certain communities to manage inflation expectations effectively.

    \item \textbf{Media as a Source of Correlated Noise:} We contribute to the literature on dispersed information and noisy signals by pointing to media narratives as a source of correlated noise. Macro models often assume exogenous noise in information or dispersed signals (e.g., \citep{Lorenzoni2009}). We give that microfoundation: media echo chambers are a plausible generator of such noise, and importantly, the noise is \emph{correlated among subsets of agents}. This can help explain why groups of agents make similar forecast errors (e.g., all underestimating inflation) in a way independent of fundamentals. This mechanism can be built into macro models of expectations (introducing a narrative noise term that hits certain fractions of agents).

    \item \textbf{Shift–Share Design in Behavioral Context:} Methodologically, we implement a novel shift–share IV strategy linking text-based media data with survey expectations. This cross-pollination of media studies and macro econometrics could be useful for other researchers. We effectively create an instrument for “exogenous sentiment shocks” in different regions using media data – something that could be extended to study, say, how news sentiment shocks affect consumption or investment. Our approach provides a template for how to use modern big data (like internet trends and transcripts) to tackle identification in behavioral questions.

    \item \textbf{Implications for Inequality:} By connecting narrative-driven expectation errors to inequality-relevant decisions, we explore a new channel for inequality dynamics. If our hypothesis H5 finds support – say we find that high-exposure low-SES households reduce spending or exit stock markets relative to others – then we've identified a feedback loop: media narratives not only reflect inequality (different groups follow different news), but potentially exacerbate it (via decisions influenced by biased beliefs). This is a contribution to the narrative of how economic inequality can be reinforced by information inequality. It could open the door for policy discussions on ensuring more balanced information reach to different socio-economic strata.

    \item \textbf{Policy and Practical Insights:} While primarily academic, our findings could have practical implications. For example, central bank communication strategies might consider engaging with or monitoring popular economic podcasts, especially ones that cater to groups whose expectations tend to be less anchored. If a certain narrative is causing expected inflation to drift, addressing that narrative (through fact-checks or targeted info campaigns) might mitigate unwarranted swings in expectations. Also, media literacy efforts could be important: if people realize they're in an echo chamber, they might seek alternative perspectives, reducing $\theta$. We will discuss such implications in the conclusion of the paper.
\end{itemize}

\paragraph{Extensions:} This project can be extended in various directions. A few promising ones:

\begin{itemize}
    \item \textbf{Other Media and Other Countries:} While we focus on U.S. podcasts, the method could apply to other media (e.g., YouTube channels, cable news shows) and other countries where surveys and media data exist. We could, for instance, study the U.K. or Eurozone if similar data can be gathered (some European countries have consumer expectation surveys and one could use social media as measure of narrative).
    \item \textbf{Expectations of Firms or Investors:} Do echo chambers affect not just households but also firms' expectations (say small business owners listening to partisan radio) or investor behavior (reddit forums creating narratives about markets)? That could tie into asset pricing anomalies. Our framework could be adapted to see if narrative shocks lead to market volatility or mispricing.
    \item \textbf{Narrative Interventions:} A more experimental extension would be to actively inject a narrative or provide an information counter-narrative to a randomly selected group and see how expectations shift. This would directly test narrative persuasion. This might be done via online surveys (e.g., show a news clip with a certain slant to some respondents).
    \item \textbf{Welfare and Policy Analysis:} If diagnostic distortions are significant, one could calibrate a model to gauge welfare loss (are decisions materially worse?). One could simulate, say, what if people had rational expectations vs what happens with narrative-induced expectations – who loses out? Possibly low-SES with wrong beliefs suffer more. This could justify certain information interventions.
    \item \textbf{Dynamic Effects and Persistence:} We plan to check if the effect of a narrative shock on expectations is persistent or fades as reality unfolds. If narratives cause overshooting but then people correct (like a boom of pessimism then mean reversion), that's different than permanently anchoring expectations to a biased view. Studying the time dynamics (with panel data tracking same person's expectations after shock dissipates) could be insightful.
\end{itemize}

In conclusion, this project will provide a comprehensive investigation into how \emph{echo chambers} in economic news impact expectation formation, utilizing state-of-the-art identification and rich data. The results will contribute to both macroeconomic theory (by highlighting the role of information structure in expectation deviations) and to practical understanding of the interplay between media, behavior, and inequality in the macroeconomy.

\vfill \pagebreak

\bibliography{references}

\end{document}


